\documentclass[11pt]{article}
\usepackage[margin=1in]{geometry}
\usepackage{amsmath,amssymb}
\usepackage{graphicx}
\usepackage{booktabs}
\usepackage{hyperref}
\usepackage{natbib}
\usepackage{setspace}
\usepackage{caption}
\usepackage{multirow}

\title{Female-Specific Cortical Circuit Hyperexcitability Driven by mGluR5--ER$\alpha$ Signaling in a \textit{Pten} Deletion Model of Autism}
\author{Author A$^{1}$, Author B$^{1}$, Author C$^{2}$, Author D$^{1}$ \\
\small $^{1}$Department of Neuroscience, University of Texas Southwestern Medical Center \\
\small $^{2}$Department of Neurology}
\date{}

\begin{document}
\maketitle

\begin{abstract}
Autism spectrum disorder (ASD) is diagnosed approximately four times more frequently in males, yet females with ASD often exhibit more severe sensory processing deficits and higher rates of epilepsy. The biological mechanisms underlying these sex differences remain poorly understood. Here we identify a female-specific cortical circuit hyperexcitability driven by mGluR5--ER$\alpha$ signaling in a \textit{Pten} deletion mouse model. Conditional deletion of \textit{Pten} in neocortical pyramidal neurons (NSEPten KO) produced a 56\% increase in UP state duration selectively in female mice, with significant sex $\times$ genotype interactions in layers 2/3 ($F_{1,73} = 5.555$, $p < 0.05$) and 5 ($F_{1,71} = 7.332$, $p < 0.01$). The mGluR5 negative allosteric modulator MTEP corrected prolonged UP states only in females, and mGluR5 protein levels were elevated in female layer-5-enriched cortical lysates. Two selective ER$\alpha$ antagonists (MPP, GNE) rescued UP state duration in female Cre(+) mice. Genetic reduction of ER$\alpha$ (Esr1$^{fl/+}$) in \textit{Pten} KO neurons rescued circuit excitability, protein synthesis, and soma size selectively in females. Female NSEPten KO mice also showed deficits in temporal sensory processing, reduced social interaction, hyperactivity, enhanced fear learning, and mGluR5-dependent seizures. The female-specific cortical hyperexcitability and seizure phenotype were replicated in the clinically relevant germline \textit{Pten}$^{+/-}$ model. These results identify an mGluR5--ER$\alpha$--ERK--protein synthesis signaling axis that produces female-selective circuit dysfunction downstream of an ASD-risk gene.
\end{abstract}

\section{Introduction}

Autism spectrum disorder (ASD) shows a marked sex bias in prevalence, with males diagnosed approximately four times more often than females \citep{loomes2017}. However, accumulating evidence suggests that ASD may manifest differently in females and is likely underdiagnosed \citep{werling2013}. Intriguingly, among individuals diagnosed with ASD, females frequently present with more severe symptoms in certain domains, including sensory processing abnormalities and higher rates of epilepsy. Understanding the biological basis of these sex differences is critical for developing sex-informed therapeutic strategies.

Mutations in the phosphatase and tensin homolog gene (\textit{PTEN}) are strongly associated with ASD, intellectual disability, macrocephaly, and epilepsy within the PTEN hamartoma tumor syndrome spectrum \citep{butler2005}. PTEN is a lipid phosphatase that negatively regulates the PI3K/Akt/mTOR signaling cascade. In the context of cancer biology, PTEN loss-of-function interacts with estrogen receptor alpha (ER$\alpha$) to drive tumor growth, but whether analogous sex-dependent mechanisms operate in the brain has been unknown.

Mouse models of \textit{Pten} deficiency have revealed social behavior deficits and neuronal hypertrophy \citep{kwon2006, page2009}. Notably, germline \textit{Pten} heterozygous mice show female-specific social behavior deficits, yet the underlying neural circuit physiology has not been explored. Meanwhile, the metabotropic glutamate receptor mGluR5 has been implicated in circuit hyperexcitability in other models of ASD and intellectual disability, including Fragile X syndrome \citep{huber2002, hays2011}. Whether mGluR5 contributes to sex-specific circuit dysfunction in \textit{Pten}-associated ASD is unknown.

Here we use extracellular recordings of spontaneous UP states in acute thalamocortical slices to investigate sex-specific effects of \textit{Pten} deletion on neocortical circuit function. UP states provide a sensitive readout of intact recurrent excitatory-inhibitory circuit dynamics and have been used to identify circuit-level phenotypes in multiple ASD models \citep{gibson2008}. We combine pharmacological, genetic, biochemical, and behavioral approaches to identify the molecular pathway responsible for female-selective hyperexcitability.

\section{Methods}

\subsection{Animal Models}

All experiments used mice on a C57BL/6 background maintained on a phytoestrogen-reduced diet (Teklad Global 16\% Protein Rodent Diet). The primary conditional knockout model (NSEPten KO) was generated by crossing NSE-Cre mice with \textit{Pten}$^{fl/fl}$ mice. NSE-Cre is expressed postnatally (beginning approximately P2) in cortical layers III--V neurons and hippocampal CA3/dentate gyrus. Greater than 99\% of targeted neurons are excitatory based on GABA immunostaining. Controls were Cre(-) littermates (\textit{Pten}$^{fl/fl}$ or \textit{Pten}$^{fl/+}$). For ER$\alpha$ genetic reduction experiments, NSE-Cre/\textit{Pten}$^{fl/fl}$/\textit{Esr1}$^{fl/+}$ mice were generated. The clinically relevant germline heterozygous model (\textit{Pten}$^{+/-}$) was produced by crossing CMV-Cre with \textit{Pten}$^{fl}$ mice (Table~\ref{tab:models}).

\begin{table}[htbp]
\centering
\caption{Animal models and genotypes.}
\label{tab:models}
\begin{tabular}{llll}
\toprule
Model & Genotype & Abbreviation & Background \\
\midrule
Conditional Pten KO & NSE-Cre(+)/Pten$^{fl/fl}$ & Cre(+) / NSEPten KO & C57BL/6 \\
Control & NSE-Cre(-)/Pten$^{fl/fl}$ or $^{fl/+}$ & Cre(-) & C57BL/6 \\
ER$\alpha$ reduction & NSE-Cre/Pten$^{fl/fl}$/Esr1$^{fl/+}$ & Cre(+)/Esr1$^{fl/+}$ & C57BL/6 \\
Germline heterozygous & Pten$^{+/-}$ & Pten$^{+/-}$ & C57BL/6J \\
\bottomrule
\end{tabular}
\end{table}

\subsection{UP State Recordings}

Thalamocortical slices (400~$\mu$m) were prepared from P18--25 (NSEPten KO) or P42--60 (\textit{Pten}$^{+/-}$) mice. Slices were recovered in standard artificial cerebrospinal fluid (ACSF) and then transferred to a high-activity ACSF containing elevated potassium (5~mM KCl) and reduced divalent cations (1~mM MgCl$_2$, 1~mM CaCl$_2$) at 32$^{\circ}$C to promote spontaneous UP states. Extracellular multiunit activity was recorded from somatosensory barrel cortex using 0.5~M$\Omega$ tungsten electrodes (FHC), amplified 10,000-fold, sampled at 2.5~kHz, and bandpass filtered at 500~Hz--3~kHz. UP states were detected when the signal amplitude exceeded 15$\times$ the RMS noise for a minimum duration of 200~ms, with a minimum inter-UP interval of 600~ms.

\subsection{Pharmacological Treatments}

Bath-applied pharmacological agents included: MTEP (3~$\mu$M, mGluR5 negative allosteric modulator, 2~hr incubation), MPP (1~$\mu$M, selective ER$\alpha$ antagonist, 1.5--2~hr), GNE (1~$\mu$M, selective ER$\alpha$ antagonist, 1.5--2~hr), U0126 (20~$\mu$M, MEK/ERK inhibitor, 2~hr), anisomycin (20~$\mu$M, protein synthesis inhibitor, 1~hr), and cycloheximide (60~$\mu$M, protein synthesis inhibitor, 1~hr).

\subsection{Molecular and Biochemical Analyses}

Western blots for mGluR5 protein levels were performed on layer-5-enriched cortical lysates. Co-immunoprecipitation assessed mGluR5--ER$\alpha$ protein complexes. De novo protein synthesis was measured using the SUnSET puromycin incorporation assay in individual layer~5 neurons \citep{schmidt2009}. Soma size was quantified from Nissl-stained sections.

\subsection{In Vivo Sensory Processing and Behavior}

Auditory steady-state responses with gap modulation (gap-ASSR) measured temporal sensory processing in P21 pups using EEG, quantified by inter-trial phase coherence (ITPC). Behavioral testing in young adults included three-chamber social interaction, open field locomotion, and fear conditioning. Seizure susceptibility was assessed via flurothyl exposure with survival analysis.

\subsection{Statistical Analysis}

Two-way ANOVA (sex $\times$ genotype) was used for UP state comparisons. Pharmacological rescue experiments used paired or unpaired $t$-tests as appropriate. Survival analysis (Mantel-Cox log-rank test) was used for seizure susceptibility. All comparisons used $\alpha = 0.05$ with appropriate corrections for multiple comparisons.

\section{Results}

\subsection{Female-Specific Prolongation of UP States in NSEPten KO Mice}

Extracellular multiunit recordings from acute thalamocortical slices revealed a striking sex-specific effect of \textit{Pten} deletion on cortical circuit dynamics. In layer~4 of somatosensory barrel cortex, female Cre(+) mice exhibited a 56\% increase in UP state duration compared to female Cre(-) controls, whereas male Cre(+) mice showed no change relative to male Cre(-) controls (Table~\ref{tab:upstates}). No sex difference was observed in Cre(-) controls, confirming that the effect requires \textit{Pten} deletion.

\begin{table}[htbp]
\centering
\caption{UP state duration results in NSEPten KO mice (P18--25). Sex $\times$ genotype interaction statistics from two-way ANOVA.}
\label{tab:upstates}
\begin{tabular}{lccccc}
\toprule
Layer & Female Cre(+) vs Cre(-) & Male Cre(+) vs Cre(-) & $F$ statistic & df & $p$-value \\
\midrule
L4 & +56\% (increased) & No change & -- & -- & -- \\
L2/3 & Increased & No change & 5.555 & (1, 73) & $< 0.05$ \\
L5 & Increased & No change & 7.332 & (1, 71) & $< 0.01$ \\
\bottomrule
\end{tabular}
\end{table}

The sex $\times$ genotype interaction was significant in both L2/3 ($F_{1,73} = 5.555$, $p < 0.05$) and L5 ($F_{1,71} = 7.332$, $p < 0.01$). No effect of sex or genotype was observed on UP state amplitude or frequency in any layer, indicating that the phenotype is specific to UP state duration.

\subsection{mGluR5 Mediates Female-Specific Hyperexcitability}

To identify the molecular mechanism underlying the female-specific UP state prolongation, we first tested the role of mGluR5 using the negative allosteric modulator MTEP (3~$\mu$M, 2~hr incubation). MTEP selectively corrected the prolonged UP state duration in female Cre(+) slices, restoring values to Cre(-) control levels. MTEP had no effect on UP state duration in female Cre(-), male Cre(+), or male Cre(-) slices.

Western blot analysis of layer-5-enriched cortical lysates revealed that mGluR5 protein levels were significantly elevated in female Cre(+) cortex compared to female Cre(-) controls. In contrast, male Cre(+) mice showed no increase in mGluR5 protein relative to male controls.

\subsection{ER$\alpha$ Is Required for Female-Specific Circuit Dysfunction}

Given the known interaction between PTEN and ER$\alpha$ in cancer biology, we tested whether ER$\alpha$ contributes to the female-specific circuit phenotype. Two structurally distinct selective ER$\alpha$ antagonists, MPP and GNE, each reduced UP state duration in female Cre(+) slices when applied to the bath. Neither compound affected UP states in male Cre(+) slices.

Genetic reduction of ER$\alpha$ provided the most compelling evidence. Crossing \textit{Pten}$^{fl/fl}$ mice with \textit{Esr1}$^{fl/+}$ mice in the NSE-Cre background produced a heterozygous reduction of ER$\alpha$ specifically in \textit{Pten}-deleted neurons. This manipulation rescued UP state duration, de novo protein synthesis (measured by puromycin incorporation), and soma enlargement---all selectively in females (Table~\ref{tab:eralpha}).

\begin{table}[htbp]
\centering
\caption{Genetic reduction of ER$\alpha$ rescues female-specific phenotypes.}
\label{tab:eralpha}
\begin{tabular}{lcccc}
\toprule
Genotype & Sex & UP State Duration & Protein Synthesis & Soma Size \\
\midrule
Cre(+) & Female & Prolonged & Elevated & Enlarged \\
Cre(+)/Esr1$^{fl/+}$ & Female & Rescued & Rescued & Rescued \\
Cre(+) & Male & Normal & Normal & Normal \\
Cre(+)/Esr1$^{fl/+}$ & Male & No change & No change & No change \\
\bottomrule
\end{tabular}
\end{table}

Co-immunoprecipitation experiments revealed elevated mGluR5--ER$\alpha$ protein complexes selectively in female cortex, providing a biochemical basis for the sex-specific signaling interaction.

\subsection{ERK and Protein Synthesis Mediate the Downstream Signaling Cascade}

To delineate the signaling pathway downstream of mGluR5--ER$\alpha$ complexes, we tested inhibitors of ERK and de novo protein synthesis. The MEK inhibitor U0126 (20~$\mu$M, 2~hr) corrected UP state duration in female Cre(+) slices without affecting female Cre(-) controls. Similarly, two mechanistically distinct protein synthesis inhibitors---anisomycin (20~$\mu$M, 1~hr) and cycloheximide (60~$\mu$M, 1~hr)---each rescued the female Cre(+) UP state phenotype.

Direct measurement of de novo protein synthesis using the SUnSET assay confirmed that puromycin incorporation was elevated in female Cre(+) layer~5 PTEN-knockout neurons relative to neighboring PTEN-positive neurons within the same sections. Male Cre(+) neurons showed no such elevation. These results establish a signaling cascade from mGluR5--ER$\alpha$ through ERK to de novo protein synthesis that drives the female-specific circuit hyperexcitability.

\subsection{Sensory Processing Deficits and Behavioral Phenotypes}

Female NSEPten KO mice displayed deficits in temporal processing of sensory stimuli, as measured by gap-ASSR at P21. Inter-trial phase coherence (ITPC) during gap-modulated auditory steady-state responses was reduced in both auditory and frontal cortex. While both sexes showed ITPC deficits with the gap stimulus, female Cre(+) ITPC was not significantly above baseline, indicating a more severe impairment.

Behavioral testing in young adults revealed that female NSEPten KO mice exhibited reduced social interaction in the three-chamber test, hyperactivity in the open field, and enhanced fear learning in the conditioning paradigm. Seizure susceptibility testing via flurothyl exposure showed that female NSEPten KO mice had increased seizure severity that was mGluR5-dependent, whereas male seizure susceptibility was not mGluR5-dependent.

\subsection{Replication in the Germline \textit{Pten}$^{+/-}$ Model}

To determine whether the female-specific circuit phenotype extends to a more clinically relevant model, we examined \textit{Pten}$^{+/-}$ mice at 6--8 weeks of age. UP state recordings revealed that female \textit{Pten}$^{+/-}$ mice showed significantly increased UP state duration and total time in UP state compared to wild-type littermates, while male \textit{Pten}$^{+/-}$ mice did not differ from wild-type (Table~\ref{tab:germline}).

\begin{table}[htbp]
\centering
\caption{UP state phenotypes in the germline \textit{Pten}$^{+/-}$ model (6--8 weeks old).}
\label{tab:germline}
\begin{tabular}{lcccc}
\toprule
Measure & Female Pten$^{+/-}$ & Female WT & Male Pten$^{+/-}$ & Male WT \\
\midrule
UP state duration & Increased & Baseline & No change & Baseline \\
Total time in UP state & Increased & Baseline & No change & Baseline \\
$N$ (slices) & 43--58 & 43--58 & 43--58 & 43--58 \\
$N$ (mice/sex/genotype) & 8--12 & 8--12 & 8--12 & 8--12 \\
\bottomrule
\end{tabular}
\end{table}

Survival analysis of repeated daily flurothyl-induced seizures (P90--100) demonstrated decreased survival specifically in female \textit{Pten}$^{+/-}$ mice (Mantel-Cox test), confirming the female-specific seizure vulnerability in the human disease-relevant model.

\section{Discussion}

We have identified an mGluR5--ER$\alpha$--ERK--protein synthesis signaling axis that produces female-selective neocortical circuit hyperexcitability downstream of the ASD-risk gene \textit{Pten}. This represents, to our knowledge, the first identification of a molecular mechanism for sex-specific effects of an ASD-risk gene on neural circuit function.

\subsection{A Convergent Signaling Mechanism}

The convergence of mGluR5 and ER$\alpha$ signaling in the brain parallels findings from breast cancer biology, where PTEN loss interacts with ER$\alpha$ to hyperactivate PI3K/Akt/mTOR. Our results extend this interaction to neural circuits, showing that enhanced mGluR5--ER$\alpha$ complexes in female cortex drive a cascade through ERK and de novo protein synthesis that ultimately manifests as prolonged UP states---a hallmark of circuit hyperexcitability \citep{gibson2008}.

The female specificity of this pathway likely reflects the higher baseline levels of ER$\alpha$ expression and/or activity in female neurons, creating a permissive environment for the mGluR5--ER$\alpha$ interaction to occur when PTEN is lost. Importantly, the use of a phytoestrogen-reduced diet throughout the study minimizes the possibility that dietary estrogens contribute to the observed sex differences.

\subsection{Clinical Implications}

Our findings have several clinical implications. First, they suggest that females with \textit{PTEN} mutations may be particularly susceptible to cortical hyperexcitability and seizures, consistent with clinical observations of sex-specific presentations in PTEN hamartoma syndromes. Second, the identification of mGluR5 as a mediator of female-specific hyperexcitability suggests that mGluR5-targeted therapeutics (such as negative allosteric modulators) might be particularly effective in female patients. Third, the involvement of ER$\alpha$ raises the possibility that hormonal status modulates the severity of circuit dysfunction, with implications for puberty and menstrual-cycle-related fluctuations in seizure susceptibility.

\subsection{Relationship to Other ASD Models}

The mGluR5-dependent circuit hyperexcitability we identify in female \textit{Pten} KO mice shares features with the mGluR5-dependent phenotypes in Fragile X syndrome models \citep{huber2002, hays2011}, suggesting convergent circuit-level pathology despite distinct genetic etiologies. However, an important distinction is the sex specificity: in Fragile X models, mGluR5-dependent phenotypes are typically observed in males (consistent with X-linked inheritance), whereas in our \textit{Pten} model, they are female-specific. This highlights the importance of examining both sexes in preclinical ASD research and suggests that different ASD-risk genes may produce distinct patterns of sex-differential vulnerability.

\subsection{Limitations}

Several limitations should be acknowledged. The NSE-Cre driver targets layers III--V excitatory neurons and is most abundant in sensory cortex, so the phenotype we describe may not generalize to all cortical regions or neuron types. The UP state assay, while sensitive to circuit-level changes, is performed in acute slices that lack intact thalamocortical and cortico-cortical connectivity. The ages examined (P18--25 for NSEPten KO, 6--8 weeks for \textit{Pten}$^{+/-}$) span a developmental window that may not capture the full trajectory of the phenotype.

\section{Conclusion}

We demonstrate that deletion of the ASD-risk gene \textit{Pten} in neocortical pyramidal neurons produces female-specific cortical circuit hyperexcitability mediated by an mGluR5--ER$\alpha$ signaling cascade through ERK and de novo protein synthesis. This sex-specific circuit dysfunction is accompanied by sensory processing deficits, behavioral abnormalities, and mGluR5-dependent seizure susceptibility, and is replicated in the clinically relevant germline \textit{Pten}$^{+/-}$ model. These findings provide a molecular framework for understanding sex differences in ASD-associated circuit dysfunction and identify potential sex-specific therapeutic targets.

\bibliographystyle{plainnat}
\bibliography{references}

\end{document}
